\documentclass[12pt, letterpaper]{article}

% --- PREÁMBULO ---
\usepackage[utf8]{inputenc}
\usepackage[T1]{fontenc}
\usepackage[spanish, es-noshorthands]{babel}
% Margen de 1 pulgada en todos los lados
\usepackage[letterpaper, margin=1in]{geometry}
\usepackage{graphicx}
\usepackage[table]{xcolor}
\usepackage[most]{tcolorbox}
\usepackage{enumitem}
\usepackage{amssymb}
\usepackage{amsmath}
\usepackage{multicol}
\usepackage{helvet}
\renewcommand{\familydefault}{\sfdefault}

% --- PAQUETE PARA FONDO ---
\usepackage{eso-pic}



% --- COLORES INSTITUCIONALES ---
\definecolor{valdark}{RGB}{50, 50, 50}
\definecolor{boxbg}{RGB}{245, 245, 245}

% --- COMANDO PARA CASILLAS DE RESPUESTA GUIADAS ---
\newcommand{\boxfill}[1][1.5cm]{\tcbox[on line, size=minimal, colframe=valdark!60, colback=boxbg, boxrule=0.8pt, arc=1mm, boxsep=3pt, left=2pt, right=2pt, top=2pt, bottom=2pt]{\makebox[#1]{\vphantom{M}}}}

\begin{document}
	
	% --- ENCABEZADO FORMATO OFICIAL ---
	\begin{center}
		\textbf{TALLER DE REPASO 6M-11}
	\end{center}

	
	% --- PUNTO 1 ---
	\noindent \textbf{1. Lee cada situación de la vida real y tradúcela a un número entero (pon el número con su signo $+$ o $-$ dentro de la caja). (20\%)}
	
	\vspace{0.4cm}
	\begin{multicols}{2}
		\begin{enumerate}[label=\textbf{\Alph*)}, itemsep=0.8cm]
			\item Bajé al sótano 3 del edificio: \boxfill[1.5cm]
			\item El termómetro marca $10^\circ$ sobre cero: \boxfill[1.5cm]
			\item Tengo \$ 25.000 en mi billetera: \boxfill[1.5cm]
			\item Un submarino está a 300 m de prof.: \boxfill[1.5cm]
			\item Le debo \$ 8.000 a mi hermano: \boxfill[1.5cm]
			\item Roma se fundó en el año 753 A.C.: \boxfill[1.5cm]
		\end{enumerate}
	\end{multicols}
	
	\vspace{0.6cm}
	
	% --- PUNTO 2 ---
	\noindent \textbf{2. Escribe el símbolo Mayor que ($>$), Menor que ($<$) o Igual ($=$) para comparar los siguientes números. ¡Cuidado con los negativos! (20\%)}
	
	\vspace{0.4cm}
	\begin{multicols}{3}
		\begin{enumerate}[label=\textbf{\Alph*)}, itemsep=0.8cm]
			\item $-150$ \boxfill[0.8cm] $-50$
			\item $+200$ \boxfill[0.8cm] $-900$
			\item $-8$ \boxfill[0.8cm] $0$
			\item $0$ \boxfill[0.8cm] $-12$
			\item $-1.000$ \boxfill[0.8cm] $-5$
			\item $+14$ \boxfill[0.8cm] $+14$
		\end{enumerate}
	\end{multicols}
	
	\vspace{0.6cm}
	
	% --- PUNTO 3 ---
	\noindent \textbf{3. Resuelve estas operaciones en tu cabeza aplicando la técnica de "Tengo y Debo". Escribe el resultado final (con su signo si es negativo). (20\%)}
	
	\vspace{0.4cm}
	\begin{multicols}{2}
		\begin{enumerate}[label=\textbf{\Alph*)}, itemsep=1.2cm]
			\item $- 150 + 400 =$ \boxfill[2cm]
			\item $- 350 - 250 =$ \boxfill[2cm]
			\item $700 - 850 =$ \boxfill[2cm]
			\item $- 500 + 500 =$ \boxfill[2cm]
			\item $- 200 + 600 - 100 =$ \boxfill[2cm]
			\item $2.000 - 2.500 - 300 =$ \boxfill[2cm]
		\end{enumerate}
	\end{multicols}
	\newpage
	
	
	% --- PUNTO 4 ---
	\noindent \textbf{4. Lee la situación, extrae los números para armar la operación matemática y calcula el resultado mentalmente. (20\%)}
	
	\vspace{0.4cm}
	\noindent \textbf{Situación:} Tienes una deuda de \textbf{\$ 15.000} en la tienda del colegio. Al día siguiente pagas \textbf{\$ 8.000} para bajar la cuenta. Pero más tarde fías unas papas por \textbf{\$ 3.000} más. ¿Cuál es tu saldo final?
	
	\vspace{0.8cm}
	\begin{center}
		\begin{tabular}{c}
			\textbf{Arma la Operación:} \\[0.4cm]
			\boxfill[2cm] \boxfill[2cm] \boxfill[2cm] $=$ \boxfill[2cm] \\[1.2cm]
			\textbf{Respuesta:} \\[0.4cm]
			Mi saldo es de \$ \boxfill[2cm] (Quedo debiendo).
		\end{tabular}
	\end{center}
	
	\vspace{1.5cm}
	
	% --- PUNTO 5 ---
	\noindent \textbf{5. Lee cada situación, escribe la operación completa en el primer recuadro grande y el resultado en el segundo. (20\%)}
	
	\vspace{0.6cm}
	\begin{enumerate}[label=\textbf{\Alph*)}, itemsep=1.5cm]
		\item Un submarino explora a \textbf{600} metros de profundidad (bajo el mar), desciende \textbf{200} metros más para investigar una cueva y luego sube \textbf{450} metros. ¿A qué profundidad quedó el submarino?
		
		$\rightarrow$ \textbf{Op:} \boxfill[6cm] $=$ \boxfill[1cm] \hspace{0.5cm} \textbf{Rta:} Quedó a \boxfill[1.5cm] m.
		
		\item La temperatura a las 6:00 a.m. era de \textbf{$4^\circ$C bajo cero}. En la tarde subió \textbf{$15^\circ$C} por el sol, y en la noche volvió a bajar \textbf{$8^\circ$C}. ¿Cuál fue la temperatura final en la noche?
		
		$\rightarrow$ \textbf{Op:} \boxfill[6cm] $=$ \boxfill[1cm] \hspace{0.5cm} \textbf{Rta:} Fue de \boxfill[1.5cm] $^\circ$C.
	\end{enumerate}
	
	\vspace{2cm}
	
	
\end{document}